\documentclass{leveldoc}
\DocUses{foundations,high_voltage_dome}
\DocFingerprint{9fe5c4cdb74eb451}

\begin{document}

\section{What the level asks}\label{sec:asks}

\leveldescription

One positive ion is launched along $+x$ from the left edge. A large metal sphere
sits on that line of flight and is grounded, so its surface stays at zero
potential. A charge placed near it draws an opposite image inside the metal. The
player may add up to three charges, of either sign, from the listed magnitudes.
There are no magnets.

The ion must enter the detector. There is no acceptance window, no kinetic window
and no gate. A flight ends when the ion touches a charge or the sphere, leaves the
world, or enters the detector (\xsecref{foundations}{sec:board}). The sphere's
obstacle is the metal plus a contact shell of thickness $\val{contact}$.

The reference adds one positive charge below the launch. The sphere alone leaves
the ion on a straight line, under the detector. The added charge lifts it onto
the detector's lower edge, passing close to the metal. This document uses the
dome's grounded image, with no centre charge, and integrates the flight by the
classical Runge--Kutta method.

\section{Setup}\label{sec:setup}

\setuptable

The speed of light is $c = \val{cLight}$. The ion has charge $q = \val{qPart}$, mass
$m = \val{mPart}$ and radius $0$, and it is launched with kinetic energy
$T_0 = \val{T0}$. A fixed charge has radius $R = \val{Rcharge}$. The sphere has
centre $(\val{sphereX}, \val{sphereY})$, radius $a = \val{sphereA}$ and potential
$\alpha = \val{alpha}$. The reference charge is $Q = \val{QPlay}$ at
$(\val{xPlay}, \val{yPlay})$.

The launch is at $(\val{xLaunch}, \val{yLaunch})$, straight along $+x$. Then
$\gamma_0 = \val{gamma0}$, $\abs{\vect{p}_0} = \val{p0}$ and
$v_0 = \val{v0} = \val{vOverC}\,c$ (\xeqref{foundations}{eq:launch}). The detector
is $\val{xDetLo} \le x \le \val{xDetHi}$, $\val{yDetLo} \le y \le \val{yDetHi}$.
The world is $-1 \le x \le \val{xWorld}$, $-1 \le y \le \val{yWorld}$.

\begin{figure}[ht]
\centering
\begin{tikzpicture}
\begin{axis}[axis equal image, width=0.92\linewidth,
  xmin=-1.5, xmax=32.5, ymin=-2, ymax=22.5,
  xlabel={$x$}, ylabel={$y$}, tick label style={font=\small}]
  \draw[black!40, dashed] (-1,-1) rectangle (\rawval{xWorld},\rawval{yWorld});
  \draw[black!25] (0,0) rectangle (30,20);
  \fill[green!50!black, fill opacity=0.18]
    (\rawval{xDetLo},\rawval{yDetLo}) rectangle (\rawval{xDetHi},\rawval{yDetHi});
  \draw[green!40!black]
    (\rawval{xDetLo},\rawval{yDetLo}) rectangle (\rawval{xDetHi},\rawval{yDetHi});
  \fill[blue!30, fill opacity=0.35] (\rawval{sphereX},\rawval{sphereY}) circle[radius=\rawval{sphereA}];
  \draw[blue!50!black] (\rawval{sphereX},\rawval{sphereY}) circle[radius=\rawval{sphereA}];
  \fill[red!70!black] (\rawval{xPlay},\rawval{yPlay}) circle[radius=\rawval{Rcharge}];
  \addplot[orange!80!black] table {pathbare.dat};
  \addplot[very thick, blue!60!black] table {pathref.dat};
  \draw[-{Stealth[length=5pt]}, thick] (\rawval{xLaunch},\rawval{yLaunch}) -- (\rawval{arrowX},\rawval{arrowY});
  \fill (\rawval{xLaunch},\rawval{yLaunch}) circle[radius=0.08];
\end{axis}
\end{tikzpicture}
\caption{The grounded sphere and the two flights. The circle is the sphere and the
disc is the reference charge; the green box is the detector. Blue is the reference
flight: it climbs past the sphere, close to the surface, and meets the bottom of
the detector. Orange is the same launch with that charge removed: it stays on the
launch line and leaves through the right edge of the world, below the detector.
Both paths stay outside the sphere.}
\label{fig:board}
\end{figure}

\section{Principles and model}\label{sec:model}

\begin{principles}
  \item The field of a point charge and the force on a non-polarizable particle,
    \xeqref{foundations}{eq:coulomb} and \xeqref{foundations}{eq:sphereforce}.
  \item The relativistic equation of motion \xeqref{foundations}{eq:eom}, the launch
    \xeqref{foundations}{eq:launch}, and energy in a static field
    \xeqref{foundations}{eq:energy}. The motion stays in the plane
    (\xsecref{foundations}{sec:plane}).
  \item A metal sphere is an equipotential. Its relaxation is instantaneous on the
    flight's time scale (\physref{2.6}). A grounded sphere is the dome's image
    \xeqref{high_voltage_dome}{eq:kelvin} with the centre charge omitted.
  \item The board's rules (\xsecref{foundations}{sec:board}): obstacles first, then
    the world, then the detector. The classical Runge--Kutta method and the cubic
    Hermite location of an event (\xsecref{foundations}{sec:numerics}).
  \item The radiated power \xeqref{foundations}{eq:larmor}, summed by Simpson's rule
    \xeqref{foundations}{eq:simpson}. It is not put back into the force.
\end{principles}

Outside the sphere the electric field is the Coulomb field of the real charge and
of its image, plus the image of the ion itself. There is no centre charge and no
magnetic field. The state $(\vect{x}, \vect{p})$ obeys \xeqref{foundations}{eq:eom}
with that $\vect{E}$, from the launch until the first ending event.

\begin{assumption}{Electrostatic images}\label{as:static}
The metal answers at once. That is exact at infinite $c$, and for finite $c$ it
drops the retardation across the board. The launch speed is $\val{vOverC}\,c$.
The game uses the same static images (\physref{2.6}), so this is the model being
compared, not a correction on top of it.
\end{assumption}

\begin{assumption}{No radiation reaction}\label{as:radiation}
The ion radiates \xeqref{foundations}{eq:larmor}, and that energy is not taken
from the flight. The larger fraction of the launch energy on the two flights here
is $\val{radFracMax}$.
\end{assumption}

\begin{assumption}{The ion's own image}\label{as:self}
The ion induces the grounded image \xeqref{high_voltage_dome}{eq:self}. Along
these two flights the largest sampled $\abs{q\vect{E}_\mathrm{self}}$ is
$\val{selfForce}$, a fraction $\val{selfBound}$ of $T_0$ per cell. The samples
are the step ends and nine Hermite stations on each step, not a continuous
maximum. The force is kept in the integration. The smallest clearance of the
surface is $\val{clearMin}$, against a shell of thickness $\val{contact}$ and a
sample span $\val{sampleArc}$. The closest centre distance is $\val{rNear}$.
\end{assumption}

\begin{assumption}{The fixed step}\label{as:step}
The reported step is $h = \val{H}$. Halving it estimates the largest error of an
end state as $\val{errMax}$ (\xsecref{foundations}{sec:numerics}). Ten times that
estimate lies below $10^{-9}$, so the tolerance on a state is the floor
$10^{-9}$.
\end{assumption}

\section{Derivation}\label{sec:derivation}

\subsection{The grounded image}\label{sec:image}

The image of a charge in a grounded sphere is \xeqref{high_voltage_dome}{eq:kelvin}.
On the surface the charge and its image cancel, so the potential is zero. The
dome then adds a centre charge to raise that potential (\xeqref{high_voltage_dome}{eq:centre}).
A grounded sphere is the same construction with that charge left out: the supply
holds the potential at zero, and the net charge on the metal is whatever the
images sum to. With one sphere there is nowhere further to image a charge, so
the construction is exact (\physref{2.6}).

The reference charge lies at distance $\val{sourceR}$ from the centre, so its
image is $\val{imageQ}$ at $(\val{imageX}, \val{imageY})$. That image is the
whole induced charge, $g = \val{inducedG}$. On sixteen points of the surface the
potential of this fixed set differs from $\alpha = \val{alpha}$ by at most
$\val{surfRes}$: roundoff of the cancellation, not a leftover potential.

\subsection{The ion's own image}\label{sec:self}

The ion induces the same grounded image, and the supply absorbs the change of
net charge, so no centre charge follows the ion. The image's potential at the
ion is \xeqref{high_voltage_dome}{eq:self}. The conserved energy is the dome's
\xeqref{high_voltage_dome}{eq:image-energy}. At launch the external part is
$q\phi_\mathrm{ext} = \val{qPhi0}$ and the image part is $\val{halfSelf0}$. The
image of the ion at launch is $\val{selfQ}$ at $(\val{selfX}, \val{selfY})$.

\subsection{The flight}\label{sec:flight}

There is no closed orbit: the reference charge and its image are two fixed
centres, and the ion's image moves. The state is advanced by
\xeqref{foundations}{eq:rk4} with the fixed step $h = \val{H}$. After each step
the cubic Hermite interpolant of the two ends is sampled at sixty-four fractions
of the step. The first sample that has touched an obstacle, left the world or
entered the detector is refined by bisection, eighty halvings
(\xsecref{foundations}{sec:numerics}). Obstacles are tested before the world, and
the world before the detector. A face coordinate is assigned when the residual
is below $10^{-6}$.

The error of the fine step is $\abs{y_h - y_{2h}}/15$. The same rule estimates
the Simpson error of the radiated energy. Both flights give the same outcome at
$h$ and at $2h$.

\section{The reference solution by hand}\label{sec:hand}

\begin{procedure}
  \item\label{st:launch} \emph{The launch} \xeqref{foundations}{eq:launch}.
    $\gamma_0 = \val{gamma0}$, $\abs{\vect{p}_0} = \val{p0}$, so
    $\vect{p}_0 = (\val{p0}, 0)$ and $v_0 = \val{v0}$.
  \item\label{st:image} \emph{The image} \xeqref{high_voltage_dome}{eq:kelvin}.
    The sphere is grounded, so there is no centre charge. The reference charge
    sits at distance $\val{sourceR}$, and its image is $\val{imageQ}$ at
    $(\val{imageX}, \val{imageY})$. The ground supplies that charge.
  \item\label{st:march} \emph{The march.} Advance
    \xeqref{high_voltage_dome}{eq:image-energy} and \xeqref{foundations}{eq:eom}
    by \xeqref{foundations}{eq:rk4} with step $\val{H}$, and locate the ending
    face as in \S\ref{sec:flight}. The state every five time units:
    \begin{center}\small\val{states}\end{center}
  \item\label{st:end} \emph{The arrival.} The ion meets the \val{RefFace} face of
    the detector at $(\val{RefX}, \val{RefY})$, at $t = \val{RefT}$, with
    momentum $(\val{RefPx}, \val{RefPy})$, $\val{RefAng}^\circ$ above $+x$.
    The ordinate is the detector's lower edge. The run takes $\val{RefSteps}$
    steps. The path only climbs: its lowest point is the launch.
  \item\label{st:check} \emph{The checks.} The potential on the sphere is off by
    at most $\val{surfRes}$. Along the flight, $W$ drifts by at most
    $\val{RefWdev}$. The step-halving estimate of the end state is
    $\val{RefErr}$. The closest centre distance is $\val{RefRclose}$, so the
    clearance of the surface is $\val{RefClear}$, outside the shell
    $\val{contact}$.
  \item\label{st:rad} \emph{The radiation.} Simpson's rule on steps of $h/2$
    gives $\val{RefRad}$, a fraction $\val{RefRadFrac}$ of $T_0$. Halving that
    step again changes the integral by the relative amount $\val{RefRadErr}$.
\end{procedure}

\begin{result}[The reference flight]
The ion enters the detector through its lower edge at
$(\val{RefX}, \val{RefY})$ at $t = \val{RefT}$, with momentum
$(\val{RefPx}, \val{RefPy})$. The level is solved.
\end{result}

The march is a few thousand steps of four evaluations. Each evaluation is the
Coulomb sum of three charges: the reference charge, its image and the ion's image.

\section{Comparison with the game}\label{sec:comparison}

\comparisontable

Nine comparisons are in the table: the reference outcome, time and end position,
the radiated energy of the reference, and the outcome, time and end position of
the grounded sphere with no added charge. The tolerance on a state is
$\val{tolFlight}$. Ten times the largest step-halving estimate $\val{errMax}$
lies below $10^{-9}$, so the tolerance is that floor. The radiation tolerance is
$10^{-4}$, relative. The state tolerance and the radiation tolerance were fixed
from these estimates before the comparison.

The largest absolute difference is the reference time, $\val{cmpRefTDiff}$. The
next is the bare ordinate, $\val{cmpBareYDiff}$. Our truncation $\val{errMax}$ is
smaller than either. Both differences sit under the verified-flight accuracy of
about $10^{-10}$ (\xsecref{foundations}{sec:verification}) and under the
tolerance.

The reference ordinate and the bare abscissa agree exactly: each is a face the
flight meets, the detector's lower edge and the world's right edge. The Simpson
estimate of the radiated energy is $\val{RefRadErr}$. The relative difference
from the game is $\val{cmpRefRadDiff}$, larger than that estimate and far under
the radiation tolerance. The two quadratures are not the same rule; the formula is.

\section{Other solutions}\label{sec:others}

The grounded sphere alone, which is the empty placement in
\code{placements.json}, does not solve the level.

The ion leaves along $+x$ and the self-image does not turn it by a visible
amount. The flight leaves through the \val{BareFace} edge of the world at
$(\val{BareX}, \val{BareY})$, at $t = \val{BareT}$, at an angle $\val{BareAng}$
degrees.
The ordinate stays on the launch line, whose height $\val{yLaunch}$ is below the
detector. The clearance of the sphere is $\val{BareClear}$, and the energy drifts
by $\val{BareWdev}$.

The reference charge is positive and sits below the launch, so it repels the ion
upward. Its image is opposite and lies inside the sphere, so the ion is also
drawn toward the metal. The climb clears the surface by $\val{RefClear}$ and
meets the detector still rising.

\section{What the game computes}\label{sec:engine}

The engine builds the same images (\physref{2.6}). For several spheres it images
again, to a finite depth, and fits a smooth remainder by fundamental solutions
inside each sphere. This level has one sphere, and it is grounded, so the second
imaging never starts and there is no centre charge. The surface residual measured
in \S\ref{sec:image} is the roundoff $\val{surfRes}$. The ion's own image is
included, and the energy diagnostic is \xeqref{high_voltage_dome}{eq:image-energy}.
The obstacle is the sphere plus the contact shell $\val{contact}$.

The flight itself is the adaptive eighth-order Runge--Kutta method of
\physref{5} and \physref{6}, verified as in \physref{7}. A verified end state is
good to about $10^{-10}$. The differences in \S\ref{sec:comparison} are inside
that. Radiation is integrated and not applied to the momentum
(\xsecref{foundations}{sec:radiation}, \physref{8}).

\section{Further reading}\label{sec:reading}

The image of a point charge in a grounded sphere, and the centre charge that
raises a sphere to a given potential, are derived for the dome
(\xeqref{high_voltage_dome}{eq:kelvin}, \xeqref{high_voltage_dome}{eq:centre}).
This level is that image with the centre charge omitted. The flight uses the
Foundations' Runge--Kutta method, Hermite event location and Simpson rule
(\xsecref{foundations}{sec:numerics}). The engine's conductor model is
\physref{2.6}.

\levelfileappendix
\end{document}
